\section{A Quantitative Quotient of the History Tree}\label{sec:quotient}
\subsection{Conditional responses and cap reflection}
For an occurrence $h$ of $v$, let $H_h(\eta)$ be the maximum conditional reward minus $\eta Y_h$, imposing all subtree caps, but not an exact payment equality at $h$. Let $H_h^{\rm pre}(\eta)$ remove only the cap at $h$, retaining every successor cap. Conditional rewards use weights $w_j/w_h$. Define $R_h(\eta)$ as the payment of the unique maximizer of $H_h(\eta)$; use $d_h(\eta)$ for the pre-cap payment. These definitions separate a removed current cap from the constrained continuation problems below it.

\begin{theorem}[Occurrence-to-vertex quotient and finite response closure]\label{thm:quotient}
Under the assumptions of Section~\ref{sec:model}, occurrences $h,h'$ of a public vertex $v$ with the same incoming price $\eta$ have identical normalized optimal continuation plans. In particular, $H_h=H_{h'}=:H_v$, $H_h^{\rm pre}=H_{h'}^{\rm pre}=:H_v^{\rm pre}$, and $R_h=R_{h'}=:R_v$. The following backward construction computes the common payment responses and optimal policies:
\begin{align}
 x_v^0(\eta)&=\left[(r_v-a_v\eta)/q_v\right]_{[l_v,h_v]},\label{eq:local}\\
 d_v(\eta)&=a_vx_v^0(\eta)+\beta\sum_wp_{vw}R_w(\eta),\label{eq:pre}\\
 R_v(\eta)&=\min\{B_v,d_v(\eta)\}.\label{eq:capped}
\end{align}
Write $d_v^-=a_vh_v+\beta\sum_wp_{vw}U_w$ for the finite constant on the left tail of $d_v$. If $B_v\geq d_v^-$, assign no barrier at $v$. Otherwise let $\alpha_v$ be the leftmost finite solution of $d_v(\alpha_v)=B_v$. The optimal local action and the incoming price sent to every successor are determined by
\begin{equation}\label{eq:reflection}
 s_v=\max\{\eta,\alpha_v\},\qquad x_v=x_v^0(s_v),
\end{equation}
where a missing barrier leaves $s_v=\eta$.

Every $R_v$ is continuous, nonincreasing, piecewise affine with constant tails and range $[L_v,U_v]$. The global universe of distinct response knots has size at most $3N$. Complete response tables can be constructed in
\begin{equation}\label{eq:arithmetic}
 A=O\bigl((N+M)N\log(N+1)\bigr)
\end{equation}
arithmetic operations and comparisons, using $O(N^2+M)$ scalar storage. For a fixed root promise, choose any finite $\eta_0$ with $R_0(\eta_0)=b_0$. Execution uses at most $N+1$ distinct incoming-price labels, drawn from $\eta_0$ and the finite cap barriers. The construction does not enumerate $H$ occurrences.
\end{theorem}
\begin{proof}
The subtree below an occurrence is determined by its displayed vertex and subsequent public transitions. The map replacing the initial occurrence by another occurrence of the same vertex preserves all suffix paths, relative probabilities, discounts, local rewards, bounds, and caps. It is consequently a bijection of the normalized feasible problems at any common $\eta$. Strict concavity makes their maximizing action vectors agree. This proves the quotient independently of the algorithm.

Remove the current cap. Separability gives \eqref{eq:local} and independent successor problems at the same incoming price, proving \eqref{eq:pre}. Their unique maximizers have continuous payment responses. Backward induction starts with a clipped affine local payment and preserves continuity, monotonicity, and affine pieces under weighted summation. Its two tails are the extreme attainable payments before the current cap.

When the current cap is redundant, this pre-cap maximizer is feasible. Otherwise feasibility implies that $B_v$ belongs to the range of $d_v$. The clipped finite response reaches its tails at finite prices, so a finite leftmost crossing exists. For $\eta<\alpha_v$, the pre-cap optimum at $\alpha_v$ has payment $B_v$. Let $z^*$ be that optimum and $z$ any policy satisfying the current cap. Writing $G$ for conditional reward and $Y$ for payment gives
\[
 G(z)-\eta Y(z)\leq G(z^*)-\alpha_v B_v+(\alpha_v-\eta)Y(z)
 \leq G(z^*)-\eta B_v.
\]
For $\eta\geq\alpha_v$, the pre-cap optimizer already satisfies the cap. These two cases prove reflection, \eqref{eq:capped}, and its implementing policy. If the payment is constant between two prices, the two optimality inequalities imply that either unique optimizer is also optimal at the other price. Thus the policy is unchanged on an equality plateau even when the supporting dual price is not unique. This also covers a cap at the minimum feasible payment and fixed local intervals.

Local clipping contributes at most two knots. Summing child responses introduces no knot outside their union. Capping can remove knots and introduces at most one crossing. Inducting over distinct \emph{vertices}, rather than occurrences, gives a global union of at most $3N$. Each of the $N$ responses can still have order $N$ segments, so total response storage is quadratic, not linear. Directly merging weighted coefficient changes at sorted knots gives \eqref{eq:arithmetic}; the graph takes an additional $O(N+M)$ records.

Finally, each executed price is the maximum of the root price and the barriers along the realized path. There are at most $N+1$ such numerical labels for this fixed root query, and at most $N(N+1)$ reachable vertex--price pairs. The root maximizer at $\eta_0$ has precisely $b_0$, so it also maximizes reward among policies satisfying the equality.\end{proof}

The unnormalized response at occurrence $h$ in Proposition~\ref{prop:laminar} is $w_hR_{v(h)}(\eta)$. This identity is the explicit link to an occurrence-indexed tree reduction. A normalized memoized implementation enjoys the same quotient bound; it is not an algorithmic competitor that our implementation is claimed to dominate.

\begin{corollary}[Several coordinates under one scalar commitment]\label{cor:coordinates}
With $d_v$ separable local coordinates, positive payment coefficients, positive quadratic coefficients, and a product interval, let $D=\sum_vd_v$. The knot universe has size at most $2D+N$, response storage is $O(N(2D+N)+M)$, and the incoming-price alphabet still has at most $N+1$ elements for a fixed query.
\end{corollary}
\begin{proof}
There are two clipping knots per coordinate, but only one scalar cap barrier per vertex. Clipping knots describe how a local action responds to a given price; they are not additional prices written into the running record. The common commitment price and cap reflection are unchanged.\end{proof}
Independent commitments can be separated only when rewards and feasible sets also separate. Coupled vector promises, nonseparable quadratic terms, and intertemporal switching do not satisfy the scalar reflection argument. Their retained primal promise-state and convex formulations remain applicable.

\subsection{Rational bit complexity}\label{sec:bit}
An arithmetic count alone does not establish polynomial time on binary inputs. Let $L$ be the total binary encoding length of all rational primitives, graph indices, and the rational root query, with rationals in lowest terms. Bit bounds below concern the coefficient-event construction, which sums affine coefficients directly. A newly computed cap crossing is used to locate an interval, \emph{not} as an input to coefficient interpolation.

\begin{theorem}[Polynomial rational preprocessing]\label{thm:bit}
All response slopes, intercepts, finite knots, cap barriers, and a returned finite root price can be represented with
\begin{equation}\label{eq:bitheight}
 S=O\bigl(NL+\log(N+M+1)\bigr)
\end{equation}
bits per numerator and denominator. The construction in Theorem~\ref{thm:quotient}, followed by one root inversion, admits the conservative bit-operation bound $O(AS^3)$ and output storage $O((N^2+M)S)$ bits. These are polynomial bounds for rational inputs, not a strongly polynomial claim or a claim that coefficient length is uniformly bounded in the horizon.
\end{theorem}
\begin{proof}
Fix an open interval between global knots. At a vertex whose current cap binds on that interval, the payment response is the constant $B_v$. Otherwise its coefficients are a local affine payment plus the discounted, probability-weighted coefficients at successors. The possible nonzero local sources are
\begin{equation}\label{eq:sources}
 B_v,\quad a_vl_v,\quad a_vh_v,\quad a_vr_v/q_v,\quad -a_v^2/q_v.
\end{equation}
Expanding either coefficient through the acyclic graph therefore gives a sum over directed paths of one source multiplied by edge probabilities and a power of $\beta$. Expansion stops at a capped vertex. No path repeats an edge, and its length is at most $N-1$.

Let $D_0$ be the product of all positive denominators and the absolute values of all nonzero numerators appearing in the rational primitive input. Its logarithm is $O(L)$. Every path-term denominator divides $D_0^{N+4}$: this exponent covers repeated discounting, each edge probability, and the products and divisions in \eqref{eq:sources}. Hence a common denominator has $O(NL)$ bits even though the number of paths may be exponential. The sum of nonnegative discounted path weights to source occurrences is at most $N$, because it is bounded by the expected number of visited vertices. The absolute value of a source is at most $2^{O(L)}$. Thus the magnitude of a coefficient is at most $N2^{O(L)}$; multiplying by the common denominator bounds its numerator by \eqref{eq:bitheight}.

Direct merging uses sums and differences of such adjacent coefficients. Its intermediate partial sums add at most a polynomial factor in $N+M$ to the preceding magnitude bound and have the same kind of path-product denominator. A new cap crossing has form $(B_v-c)/m$ on a nonconstant line $m\eta+c$; local clipping knots and a root inverse are ratios of a constant number of the same bounded-height rationals. Consequently their heights also satisfy \eqref{eq:bitheight}. Exact comparisons use cross products of bounded-height integers. Elementary integer arithmetic and reduction by greatest common divisors cost at most $O(S^3)$ per rational operation under a conservative implementation. Multiplying by \eqref{eq:arithmetic} proves the stated time bound; counting stored scalars proves the storage bound.\end{proof}

The theorem bounds preprocessing and inversion, not the decimal length of a complete printed primal--dual certificate by the same $S$. Different realized barriers can introduce different denominators into a sum of rewards. A straightforward polynomial bound for aggregate execution values is larger; the companion gives one and the experiments separately time certificate generation and audit. This distinction is practically relevant even when each response coefficient is modest.

\subsection{A compact exact certificate}\label{sec:certificate}
For each reached $(v,\eta)$, propagate discounted occupation weight $w$, effective price $s$, action $x$, and multipliers $\chi=s-\eta$, $\ell,u\geq0$. They satisfy
\begin{equation}\label{eq:kkt}
 r_v-q_vx-a_vs+\ell-u=0,\quad
 \ell(x-l_v)=u(h_v-x)=\chi(B_v-R_v(\eta))=0.
\end{equation}
Weight $\beta wp_{vw}$ moves to $(w,s)$. A backward conditional-payment calculation verifies every cap; forward flow verifies occupation weights and the root promise. The resulting dual upper bound is
\begin{equation}\label{eq:dual}
 \eta_0b_0+\sum_{(v,\eta)}w\left[
 \chi B_v+u h_v-\ell l_v+
 \frac{(r_v-a_vs+\ell-u)^2}{2q_v}\right].
\end{equation}
Ancestor prices telescope under the flow equations, and \eqref{eq:kkt} makes \eqref{eq:dual} equal the implemented reward. These rational checks certify a particular optimum; verifying entire response-curve identities is a separate implementation test. Neither kind of code execution replaces the mathematical proof above.
